\documentclass[twoside]{article}

\usepackage[preprint]{aistats2027}
\usepackage[round,authoryear]{natbib}
\usepackage{amsmath,amssymb,amsthm}
\usepackage{mathtools}
\usepackage{microtype}
\usepackage{booktabs}
\usepackage{url}

\renewcommand{\bibname}{References}
\renewcommand{\bibsection}{\subsubsection*{\bibname}}
\bibliographystyle{apalike}

\newtheorem{theorem}{Theorem}
\newtheorem{proposition}[theorem]{Proposition}
\newtheorem{corollary}[theorem]{Corollary}
\newtheorem{lemma}[theorem]{Lemma}
\newcommand{\E}{\mathbb E}
\newcommand{\cA}{\mathcal A}
\newcommand{\cR}{\mathcal R}
\newcommand{\1}{\mathbf 1}

\begin{document}
\runningauthor{}

\twocolumn[
\aistatstitle{Supplement: Characterizing Perfect Truthfulness in Continuous Online Calibration}

\aistatsauthor{Pahan Dewasurendra}

\aistatsaddress{Johns Hopkins University}
]

\appendix

\section{Probability Trees and Proper Scores}
\label{app:notation}

Let $\cA=[k]$ and let $\Delta_k$ be the probability simplex. For depth
$d$, define the report-tree space
\begin{equation*}
 \cR_d=\prod_{t=1}^d\prod_{h\in\cA^{t-1}}\Delta_k.
\end{equation*}
We write $r_h\in\Delta_k$ for the report at history $h$. The report tree
$R\in\cR_d$ induces a law $Q_R$ on $\cA^d$ by
\begin{equation}
 Q_R(x)=\prod_{t=1}^d r_{x_{<t},x_t}.
 \label{eq:supp-tree-law}
\end{equation}
On the relative interior $\cR_d^\circ$, this is a bijection between report
trees and full-support laws on $\cA^d$. Its inverse takes the usual
conditional probabilities. Both maps are continuous on the relative
interiors.

A realized-path score has the form
\begin{equation}
 S(x,R)=C(x,(r_{x_{<t}})_{t=1}^d).
 \label{eq:supp-path-score}
\end{equation}
Perfect truthfulness means that $S$ is a proper score for the joint law:
for every law $P$ and every report tree $R$,
\begin{equation}
 \E_P S(X,R_P)\leq \E_P S(X,R),
 \label{eq:supp-joint-proper}
\end{equation}
where $R_P$ is any conditional-probability tree representing $P$. Reports
at histories outside the support of $P$ do not affect either side. On
full-support laws, $R_P$ is unique.

We use losses, so lower is better. For a finite outcome space $\Omega$, a
proper loss is a continuous vector map $s:\Delta(\Omega)\to\mathbb R^\Omega$
such that
\begin{equation}
 p\cdot s(p)\leq p\cdot s(q)
 \quad\text{for all }p,q\in\Delta(\Omega).
 \label{eq:supp-proper}
\end{equation}
Its entropy or Bayes risk is $H_s(p)=p\cdot s(p)$.

\section{Continuous Proper-Score Lemmas}
\label{app:convex}

For completeness, we prove the two convex-analytic lemmas used in the main
paper. They hold on any finite outcome space and require only continuity.

\begin{lemma}[Continuous support]
\label{lem:supp-continuous}
Let $s:\Delta(\Omega)\to\mathbb R^\Omega$ be continuous and proper. For
every relative-interior $p$ and every tangent vector $v$ satisfying
$\1\cdot v=0$,
\begin{equation}
 D H_s(p)[v]=v\cdot s(p).
 \label{eq:supp-directional}
\end{equation}
The derivative is continuous in $p$.
\end{lemma}

\begin{proof}
Choose $t>0$ small enough that $p+tv$ remains in the relative interior.
Propriety at truth $p+tv$ with report $p$ gives
\begin{equation*}
 H_s(p+tv)\leq(p+tv)\cdot s(p)
 =H_s(p)+t v\cdot s(p).
\end{equation*}
Propriety at truth $p$ with report $p+tv$ gives
\begin{align*}
 H_s(p)&\leq p\cdot s(p+tv)\\
 &=H_s(p+tv)-t v\cdot s(p+tv).
\end{align*}
Therefore
\begin{equation}
 v\cdot s(p+tv)\leq
 \frac{H_s(p+tv)-H_s(p)}{t}
 \leq v\cdot s(p).
 \label{eq:supp-squeeze}
\end{equation}
Letting $t\downarrow0$ and using continuity proves the right directional
derivative. Replacing $v$ by $-v$ proves the left derivative. The expression
$v\cdot s(p)$ is continuous in $p$.
\end{proof}

\begin{corollary}[Entropy uniqueness]
\label{cor:supp-unique}
Let $s$ and $u$ be continuous proper losses on a finite outcome space. If
$H_s(p)-H_u(p)=c$ for every relative-interior $p$, then
\begin{equation}
 s_i(p)-u_i(p)=c
 \quad\text{for every }i,p.
 \label{eq:supp-unique}
\end{equation}
The conclusion extends to the simplex boundary.
The componentwise conclusion on the relative interior requires only that the
scores be continuous and proper when truths and reports are restricted to the
relative interior.
\end{corollary}

\begin{proof}
Lemma~\ref{lem:supp-continuous} gives
$v\cdot(s(p)-u(p))=0$ for every tangent vector $v$. Hence
$s(p)-u(p)=\lambda(p)\1$. Taking the inner product with $p$ shows
$\lambda(p)=c$. Continuity extends the identity to the boundary.
\end{proof}

\begin{lemma}[Shared component]
\label{lem:supp-shared}
Let $s$ and $u$ be continuous proper losses on $\Delta_k$. Fix $i\in[k]$.
If
\begin{equation*}
 s_j(q)=u_j(q)\quad\text{for every }j\ne i\text{ and }q\in\Delta_k,
\end{equation*}
then $s_i(q)-u_i(q)$ is constant in $q$.
\end{lemma}

\begin{proof}
Work first in the relative interior. Set
$d(q)=s_i(q)-u_i(q)$ and
\begin{equation*}
 F(q)=H_s(q)-H_u(q)=q_i d(q).
\end{equation*}
By Lemma~\ref{lem:supp-continuous}, for every tangent vector $v$,
\begin{equation}
 D F(q)[v]=v\cdot(s(q)-u(q))=v_i d(q).
 \label{eq:supp-shared-gradient}
\end{equation}
If $v_i=0$, the derivative is zero. Each relative-interior slice with fixed
$q_i=z$ is connected, so $F(q)=f(z)$ for some continuously differentiable
function $f:(0,1)\to\mathbb R$.

Choose any $j\ne i$ and move in direction $e_i-e_j$. Equation
\eqref{eq:supp-shared-gradient} gives
\begin{equation*}
 f'(z)=d(q)=\frac{f(z)}{z}.
\end{equation*}
Thus $(f(z)/z)'=0$, so $f(z)=cz$ and $d(q)=c$. Continuity extends the result
to the boundary.
\end{proof}

\section{Proof of the Path-Score Characterization}
\label{app:structure}

We restate the representation in report-tree notation.

\begin{theorem}[Path-score characterization]
\label{thm:supp-structure}
Let $C$ be a continuous realized-path score of depth $d$ on a finite alphabet.
It is proper for every law on $\cA^d$ if and only if there are continuous
proper one-step losses $\ell_h$ such that
\begin{equation}
 C(x,(r_{x_{<t}})_t)=\sum_{t=1}^d
 \ell_{x_{<t}}(x_t,r_{x_{<t}}).
 \label{eq:supp-additive}
\end{equation}
\end{theorem}

\begin{proof}
The reverse implication follows from iterated conditioning and one-step
propriety. We prove the forward implication by induction on $d$. The claim is
immediate for $d=1$.

Assume the claim through depth $d-1$ and consider a depth-$d$ score. We first
work with interior root reports and full-support continuation laws. Fix an
interior root report $q\in\Delta_k^\circ$. For root outcome $i$, define the
continuation score
\begin{equation}
 C_i^q(z,R)=C(iz,(q,(r_{z_{<t}})_t)),
 \label{eq:supp-cont-score}
\end{equation}
where $z\in\cA^{d-1}$ and $R\in\cR_{d-1}$. The index alignment in
\eqref{eq:supp-cont-score} is only notational: $R$ supplies every report after
the root on branch $i$.

For fixed $q$, $C_i^q$ is proper for every continuation law $P_i$. To verify
this, construct a joint law whose truthful root report is $q$, whose branch
$i$ continuation is $P_i$, and whose other continuations are arbitrary and
full support. Compare the truthful tree with a tree that changes only the
continuation reports in branch $i$. Joint propriety gives the continuation
properness inequality multiplied by $q_i>0$.

Let
\begin{equation}
 B_i(q,P_i)=\E_{Z\sim P_i}C_i^q(Z,R_{P_i})
 \label{eq:supp-branch-bayes}
\end{equation}
be the continuation Bayes risk. Fix a tuple of full-support continuation laws
$P=(P_1,\ldots,P_k)$. The vector
\begin{equation*}
 B(q,P)=(B_1(q,P_1),\ldots,B_k(q,P_k))
\end{equation*}
is a continuous proper root loss. Indeed, for a true root law $p$, compare
the fully truthful report tree with a tree that reports $q$ at the root and
reports every continuation $P_i$ truthfully. Joint propriety gives
\begin{equation*}
 \sum_i p_iB_i(p,P_i)\leq\sum_i p_iB_i(q,P_i).
\end{equation*}

Fix reference continuation laws $P_i^0$ and an interior reference root report
$q^0$. If $P_i$ changes while every other continuation stays fixed, the two
proper root losses agree in every component except $i$. By
Lemma~\ref{lem:supp-shared},
\begin{equation*}
 B_i(q,P_i)-B_i(q,P_i^0)
\end{equation*}
is independent of $q$. Define
\begin{align}
 A_i(q)&=B_i(q,P_i^0),\notag\\
 K_i(P_i)&=B_i(q^0,P_i)-B_i(q^0,P_i^0).
 \label{eq:supp-AK}
\end{align}
Then
\begin{equation}
 B_i(q,P_i)=A_i(q)+K_i(P_i).
 \label{eq:supp-bayes-separate}
\end{equation}

The vector $A(q)=(A_i(q))_{i=1}^k$ is a proper root loss. For any fixed
continuation tuple $P$, the proper loss $B(q,P)$ differs from $A(q)$ by the
outcome-dependent constant vector $(K_i(P_i))_i$. Such constants cancel from
the properness comparison at a fixed true root law.

It remains to separate the full continuation score, not only its Bayes risk.
For fixed $i$ and root reports $q,q^0$, both $C_i^q$ and $C_i^{q^0}$ are
continuous proper scores on the finite outcome space $\cA^{d-1}$. Under the
bijection in \eqref{eq:supp-tree-law}, their entropies at a full-support
continuation law $P_i$ are $B_i(q,P_i)$ and $B_i(q^0,P_i)$. By
\eqref{eq:supp-bayes-separate}, their difference is the constant
$A_i(q)-A_i(q^0)$. The interior version of
Corollary~\ref{cor:supp-unique} gives
\begin{equation}
 C_i^q(z,R)=C_i^{q^0}(z,R)+A_i(q)-A_i(q^0)
 \label{eq:supp-score-separate}
\end{equation}
for every full-support continuation report tree $R$ and every $z$.

Define
\begin{equation*}
 D_i(z,R)=C_i^{q^0}(z,R)-A_i(q^0).
\end{equation*}
Then $C_i^q(z,R)=A_i(q)+D_i(z,R)$. The score $D_i$ is continuous,
realized-path local, and proper on every depth-$(d-1)$ continuation law.
Apply the induction hypothesis to each $D_i$. Together with the proper root
loss $A$, this gives \eqref{eq:supp-additive} for interior reports.

Interior report trees are dense in $\cR_d$. Every term and the original
score are continuous, so the identity extends to all boundary reports. The
properness inequalities for the component losses also extend to boundary
true laws and reports by continuity.
\end{proof}

\section{Normalization and Risk Comparison}
\label{app:risk}

\begin{corollary}[Canonical normalization]
\label{cor:supp-normalize}
Suppose the score in Theorem~\ref{thm:supp-structure} satisfies
\begin{equation}
 C(x,(e_{x_t})_{t=1}^d)=0\quad\text{for every }x\in\cA^d.
 \label{eq:supp-endpoint}
\end{equation}
Then the components can be chosen to satisfy
\begin{equation*}
 \ell_h(i,e_i)=0,\qquad \ell_h(i,q)\geq0.
\end{equation*}
\end{corollary}

\begin{proof}
Starting from any representation, define
\begin{equation*}
 \widetilde\ell_h(i,q)=\ell_h(i,q)-\ell_h(i,e_i).
\end{equation*}
Subtracting an outcome-dependent constant preserves propriety. Propriety at
the deterministic truth $e_i$ gives
$\ell_h(i,e_i)\leq\ell_h(i,q)$, so every normalized loss is nonnegative.
Finally, for each path $x$,
\begin{equation*}
 \sum_{t=1}^d\ell_{x_{<t}}(x_t,e_{x_t})
 =C(x,(e_{x_t})_t)=0.
\end{equation*}
Thus replacing every component by its normalized version leaves $C$ exactly
unchanged.
\end{proof}

\begin{lemma}[One-step comparison]
\label{lem:supp-one-step}
Let $\ell$ be a nonnegative proper categorical loss normalized at every
simplex vertex. For interior $p,q\in\Delta_k$, define
\begin{align*}
 H(p)&=\sum_i p_i\ell(i,p),\\
 L(p,q)&=\sum_i p_i\ell(i,q).
\end{align*}
Then
\begin{equation}
 L(p,q)\leq
 \left(\max_i\frac{p_i}{q_i}\right)
 \left(\max_i\frac{q_i}{p_i}\right)H(p).
 \label{eq:supp-one-step}
\end{equation}
\end{lemma}

\begin{proof}
Nonnegativity gives
\begin{equation*}
 L(p,q)\leq\left(\max_i\frac{p_i}{q_i}\right)
 \sum_i q_i\ell(i,q).
\end{equation*}
Properness at truth $q$ gives
\begin{equation*}
 \sum_i q_i\ell(i,q)\leq\sum_i q_i\ell(i,p).
\end{equation*}
A second use of nonnegativity gives
\begin{equation*}
 \sum_i q_i\ell(i,p)
 \leq\left(\max_i\frac{q_i}{p_i}\right)H(p).
\end{equation*}
Multiplying the factors proves the claim.
\end{proof}

\begin{theorem}[Sequential comparison]
\label{thm:supp-ratio}
Let $C_T$ be continuous, perfectly truthful, and satisfy
\eqref{eq:supp-endpoint}. For interior $p,q\in\Delta_k$ and i.i.d. outcomes
$X_t\sim p$,
\begin{equation}
 \E_p C_T(X,q\1_T)\leq\kappa(p,q)\E_p C_T(X,p\1_T),
 \label{eq:supp-sequential-ratio}
\end{equation}
where
\begin{equation*}
 \kappa(p,q)=\left(\max_i\frac{p_i}{q_i}\right)
             \left(\max_i\frac{q_i}{p_i}\right).
\end{equation*}
\end{theorem}

\begin{proof}
Use the normalized components from Corollary~\ref{cor:supp-normalize}. At
each history $h$, Lemma~\ref{lem:supp-one-step} applies to $\ell_h$. Under
i.i.d. truth $p$, the conditional law of $X_t$ given $X_{<t}=h$ is $p$.
Therefore
\begin{align*}
 &\E_p[\ell_{X_{<t}}(X_t,q)\mid X_{<t}]\\
 &\qquad\leq\kappa(p,q)
 \E_p[\ell_{X_{<t}}(X_t,p)\mid X_{<t}].
\end{align*}
Take expectations and sum over $t$.
\end{proof}

For binary $p=(1-a,a)$ and $q=(1-b,b)$, if $b>a$ then
\begin{align*}
 \max_i\frac{p_i}{q_i}&=\frac{1-a}{1-b},\\
 \max_i\frac{q_i}{p_i}&=\frac{b}{a}.
\end{align*}
Their product is $b(1-a)/(a(1-b))$. The case $b<a$ gives its reciprocal.
Thus
\begin{equation}
 \kappa(p,q)=
 \exp\!\left(|\operatorname{logit}(a)-\operatorname{logit}(b)|\right).
 \label{eq:supp-binary-kappa}
\end{equation}

\section{Sharpness of the Binary Constant}
\label{app:sharpness}

We give a direct construction rather than invoke the full representation
theory of binary proper losses. Let $w:[0,1]\to\mathbb R_{\geq0}$ be a
continuous density. Define
\begin{align}
 \ell(0,q)&=\int_0^q t w(t)\,dt,\label{eq:supp-binary-loss0}\\
 \ell(1,q)&=\int_q^1 (1-t)w(t)\,dt.
 \label{eq:supp-binary-loss1}
\end{align}
The losses are continuous, nonnegative, and normalized at $q=0$ and $q=1$.
For truth $a$, differentiation of the cross-risk gives
\begin{equation*}
 \frac{d}{dq}\left((1-a)\ell(0,q)+a\ell(1,q)\right)
 =(q-a)w(q).
\end{equation*}
Hence the loss is proper.

Fix $b>a$. Choose a sequence of continuous unit-mass densities $w_n$
supported inside $(a,b)$ and concentrating at $b$ from below. Since all mass
lies above $a$ and below $b$,
\begin{align*}
 L_n(a,a)&=a\int_a^b(1-t)w_n(t)\,dt,\\
 L_n(a,b)&=(1-a)\int_a^b t w_n(t)\,dt.
\end{align*}
Therefore
\begin{equation*}
 \frac{L_n(a,b)}{L_n(a,a)}
 \longrightarrow\frac{b(1-a)}{a(1-b)}=\kappa(p,q).
\end{equation*}
For $b<a$, choose densities supported inside $(b,a)$ and concentrating at
$b$ from above. Then
\begin{align*}
 L_n(a,a)&=(1-a)\int_b^a t w_n(t)\,dt,\\
 L_n(a,b)&=a\int_b^a(1-t)w_n(t)\,dt,
\end{align*}
and the ratio approaches $a(1-b)/(b(1-a))$. This proves the sharpness claim
for every $\epsilon>0$.

\section{Batch ATB and Its Online Deviation}
\label{app:atb}

For completeness, define the empirical Averaged Two-Bin calibration error
of \citet{hartline2026perfect} by
\begin{align}
 \operatorname{ATB}_T(r,x)
 &=\frac{1}{T^2}\int_0^1\Bigg[
 \left(\sum_{t:r_t<z}(r_t-x_t)\right)^2\notag\\
 &\qquad+
 \left(\sum_{t:r_t\geq z}(r_t-x_t)\right)^2\Bigg]dz.
 \label{eq:supp-atb}
\end{align}
The choice of strict inequality at the threshold is immaterial under the
integral.

If $r=b\1_T$, all reports occupy one bin for almost every $z$. Hence
\begin{equation*}
 \operatorname{ATB}_T(b\1_T,x)
 =\frac{1}{T^2}\left(\sum_{t=1}^T(b-x_t)\right)^2.
\end{equation*}
For i.i.d. $X_t\sim\operatorname{Ber}(a)$, bias and variance give
\begin{equation}
 \E_a\operatorname{ATB}_T(b\1_T,X)
 =(b-a)^2+\frac{a(1-a)}{T}.
 \label{eq:supp-atb-expectation}
\end{equation}
Setting $b=a$ yields truthful expected ATB $a(1-a)/T$.

We next verify the two-round online deviation in the main paper. Let
$a=1/2$, set $R_1=1/2$, and set $R_2=1/4$ after $X_1=0$ or $R_2=3/4$
after $X_1=1$. Integrating \eqref{eq:supp-atb} over the intervals induced by
the two reports gives the following table.
\begin{center}
\begin{tabular}{ccc}
\toprule
$X_1X_2$ & $(R_1,R_2)$ & $\operatorname{ATB}_2$\\
\midrule
$00$ & $(1/2,1/4)$ & $1/8$\\
$01$ & $(1/2,1/4)$ & $1/16$\\
$10$ & $(1/2,3/4)$ & $1/16$\\
$11$ & $(1/2,3/4)$ & $1/8$\\
\bottomrule
\end{tabular}
\end{center}
The outcomes are equiprobable, so the adaptive expected loss is
\begin{equation*}
 \frac14\left(\frac18+\frac1{16}+\frac1{16}+\frac18\right)
 =\frac3{32}.
\end{equation*}
For truthful constant reports, \eqref{eq:supp-atb-expectation} with $T=2$
gives $1/8$. Thus the adaptive reports strictly improve on truth.

\bibliography{references}

\end{document}
